% ===================================================================================================================
% Section: Experiments (25%). Needs: amsmath, booktabs, graphicx, natbib; \graphicspath{{figures/}}; references.bib.
% Every figure and table sits in this section, next to the text that cites it.
% Sources: epills/outputs/runs/*/eval/*.json, docs/stage2_results.md, docs/evidence.md.
% ===================================================================================================================
\section{Experiments}
\label{sec:experiments}

\subsection{Setup}
\label{sec:setup}

\paragraph{Test sets.} No test medicine appears in training in any form, and every setting is fixed on validation
before a test set is scored, once. \emph{C3PI held-out}: 4,005 photos of 200 appearance groups (the medicines of
ePillID's official holdout fold and their look-alikes), photographed like the training data.
\emph{MEDISEG}~\citep{mediseg2026}: 13,973 pills of 32 Hong Kong products, taken with other phones and in cluttered
scenes; median pill size 81\,px. \emph{ePillID}~\citep{usuyama2020epillid}: the published benchmark (fold 4,
two-sided queries against references of all 4,902 types).

\paragraph{Healthcare-relevant metrics.} We simulate patients with 10 enrolled medicines (3 photos each). For the
patient's own pills we report the share \emph{named correctly}, \emph{named wrongly} (the safety-critical error),
\emph{retake} and \emph{called unknown}; for pills that are not theirs, the share \emph{rejected}. A wrong name may be
acted on, whereas ``unknown'' only costs time, so wrong names matter most. We also report top-$k$ accuracy, AUROC,
ePillID's MAP and GAP, and the reliability of the stated confidence. 95\% intervals come from resampling medicines.

\paragraph{Baselines.} (i) Run 1: the same network with a common recipe (new layers at $10^{-3}$, no margin
warm-up). (ii) The published ePillID models~\citep{usuyama2020epillid}. (iii) Look only versus look plus imprint, for
the gallery and the catalogue. (iv) For the pill finder that feeds the model, an earlier finder that returned exactly
one box on only 28\% of single-pill photos; ours counts 85--88\% of 1-, 2- and 3-pill photos right
(Figure~\ref{fig:finder}).

\begin{figure}[t]
  \centering
  \includegraphics[width=0.55\linewidth]{fig17_pill_finder.pdf}
  \caption{The pill finder on C3PI photos with known pill counts (2- and 3-pill photos composed from single-pill
    photos). Small shares are listed after each bar in legend order.}
  \label{fig:finder}
\end{figure}

\subsection{Main results}
\label{sec:main}

Table~\ref{tab:main} gives the test results and Figure~\ref{fig:main} the three headline measures. In the conditions
it was trained for, the served model (run 4) names 76.0\% of a patient's own pills correctly and only 0.04\% wrongly
(95\% CI 0.02--0.07\%), while rejecting 94.8\% of pills that are not the patient's (Figure~\ref{fig:openset}a). The
imprint check raises correct names to 78.3\% (wrong 0.06\%) by settling retakes. Against the catalogue, the right
product comes first for 80.2\% of held-out pills and is in the top 5 for 94.4\%, against 71.6\% and 91.7\% by look
alone. Adding reference photos (run 3) nearly tripled ePillID MAP@1 (13.7$\to$37.1\%), and robustness augmentation
(run 4) gained 6.4 points on MEDISEG at a small cost in C3PI conditions; run 4 is the served model.

\begin{table}[t]
  \caption{Test results on medicines never seen in training (\%; 10-medicine galleries, 3 enrolment photos,
    calibrated $\tau$ and $m$). Bottom block: the served pipeline with imprint reading (GPU reader; CPU reader in
    brackets).}
  \label{tab:main}
  \centering
  \small
  \begin{tabular}{@{}lccc@{}}
    \toprule
    & Run 2 & Run 3 & Run 4 (served) \\
    \midrule
    C3PI held-out: top-1 among 200 medicines & 82.0 & \textbf{84.4} & 83.2 \\
    \quad gallery: named correctly / wrongly & 78.6 / 0.1 & \textbf{80.4} / 0.0 & 76.0 / 0.04 \\
    \quad other medicines rejected & 94.2 & 94.1 & \textbf{94.8} \\
    MEDISEG: top-1 among 32 products & 38.9 & 41.8 & \textbf{48.2} \\
    \quad gallery: named correctly / wrongly & 17.6 / 1.1 & 18.7 / 1.0 & 18.7 / \textbf{0.7} \\
    \quad own called unknown / other rejected & 51.5 / 78.7 & 64.2 / 89.3 & 70.4 / \textbf{93.1} \\
    ePillID strict (no test type in training): MAP@1 & 13.7 & 37.1 & \textbf{42.3} \\
    \midrule
    C3PI gallery + imprint check: named correctly / wrongly & & & 78.3 (77.3) / 0.06 \\
    \quad retake / called unknown & & & 7.2 / 14.5 \\
    Catalogue (4,936 products) + imprint: top-1 / top-5 & & & 80.2 (79.2) / 94.4 \\
    \quad look only: top-1 / top-5 & & & 71.6 / 91.7 \\
    \bottomrule
  \end{tabular}
\end{table}

\begin{figure}[t]
  \centering
  \includegraphics[width=\linewidth]{fig07_main_results.pdf}
  \caption{Test results on medicines never seen in training, each scored once. Run 3 adds C3PI studio reference
    photos; run 4 adds augmentation for other phones and jittered crops.}
  \label{fig:main}
\end{figure}

\begin{figure}[t]
  \centering
  \includegraphics[width=\linewidth]{fig08_open_set_outcomes.pdf}
  \caption{What the patient sees (run 4, embedding only; 10 medicines, 3 enrolment photos). Small shares are listed
    after each bar in legend order. In C3PI conditions most own pills are named correctly and wrong names are rare;
    on other phones the threshold transfers badly and most own pills are called unknown, which is safe but
    unhelpful.}
  \label{fig:openset}
\end{figure}

\subsection{Ablations and hyperparameters}
\label{sec:ablations}

Table~\ref{tab:ablation} isolates each design decision. The learning rates matter most. With the new layers at
$10^{-3}$ the embedding collapses to about 3 effective dimensions and validation top-1 stalls at 77\%, while lower rates
keep 12--15 dimensions and reach 98\% (Figure~\ref{fig:training}). The decision settings are smooth
(Figure~\ref{fig:calib}): the crop window moves correct names by under one point, and the margin trades correct names
for retakes almost one for one. The imprint weight has a clear optimum (Figure~\ref{fig:identify}b), because a large
$\alpha$ lets false OCR reads promote wrong products; the imprint helps most when the pill's own imprint is actually
read (Figure~\ref{fig:identify}a).

\begin{table}[t]
  \caption{Ablations (validation unless marked \emph{test}).}
  \label{tab:ablation}
  \centering
  \small
  \begin{tabular}{@{}p{0.34\linewidth}p{0.46\linewidth}p{0.12\linewidth}@{}}
    \toprule
    Change & Effect & Kept \\
    \midrule
    Run 1$\to$2: lower LRs ($3{\times}10^{-5}$, $10^{-5}$), separate ArcFace LR, margin warm-up & top-1 in 10
      medicines 77.2$\to$98.3\%; AUROC 0.78$\to$0.981; correct at 95\% rejection 17$\to$90.7\%; dims
      $\approx$3$\to$12--15 & yes \\
    Run 2$\to$3: + studio reference photos & validation 98.1\%; \emph{test} ePillID MAP@1 13.7$\to$37.1\% & yes \\
    Run 3$\to$4: + photo and crop-jitter augmentation & \emph{test} MEDISEG 41.8$\to$48.2\%; C3PI 84.4$\to$83.2\%
      & yes \\
    Crop window 1.3--2.0$\times$ & named correctly 86.1--87.2\% & 1.55 \\
    Margin $m$ 0--0.3 & named correctly 87$\to$54\%; retakes 0$\to$34\% & 0.115 \\
    Imprint re-ranking $\alpha$, $K$ & top-1 72.1\% ($\alpha{=}0$)$\to$81.6\%; \emph{test} 71.6$\to$80.2\%
      & 0.05, 100 \\
    Gallery check: distinct tokens only & named wrongly 35$\to$30; other pills accepted 6,158$\to$6,150 & yes \\
    Official ePillID protocol (runs 5--7) & ePillID MAP 0.57--0.60; MEDISEG 52--54\% & benchmark \\
    \bottomrule
  \end{tabular}
\end{table}

\begin{figure}[t]
  \centering
  \includegraphics[width=\linewidth]{fig05_training_curves.pdf}
  \caption{Training. (a) Validation top-1 within a gallery of 10 medicines never trained on. (b) Effective embedding
    dimensions: with the new layers at $10^{-3}$ the embedding collapsed to about 3.}
  \label{fig:training}
\end{figure}

\begin{figure}[t]
  \centering
  \includegraphics[width=\linewidth]{fig12_identify_topk.pdf}
  \caption{Catalogue identification. (a) Top-$k$ on held-out medicines: the imprint helps most when the pill's own
    imprint is read (92.9\% first vs.\ 68.6\% when it is not). (b) Validation grid for $\alpha$ and $K$; the box marks
    the served setting.}
  \label{fig:identify}
\end{figure}

\subsection{Robustness}
\label{sec:robustness}

\emph{Other phones} are the main weakness (Table~\ref{tab:main}, Figure~\ref{fig:openset}). On MEDISEG, top-1 drops to
48\% and the C3PI-calibrated threshold calls 70\% of own pills unknown, though wrong names stay below 1\%: the system
fails safe but is not yet usable there. Small pills are the main factor (Figure~\ref{fig:robust}a: 36--37\% top-1 under
100\,px, 62\% above), and more enrolment photos help only partly (Figure~\ref{fig:robust}b). \emph{Pills outside the
catalogue}: 69.8\% of MEDISEG pills get \emph{no confident match} and 7.0\% \emph{possible match}, while 75.8\% of
catalogue pills get \emph{possible match}, and the stated confidence tracks observed accuracy
(Figure~\ref{fig:confidence}). \emph{CPU deployment}: with the smaller reader, identification loses one point (79.2\%)
and wrong names do not rise (Table~\ref{tab:main}).

\begin{figure}[t]
  \centering
  \includegraphics[width=\linewidth]{fig09_robustness.pdf}
  \caption{Robustness on MEDISEG (32 Hong Kong products, other phones, none in training). (a) Top-1 by pill size in the
    photo. (b) MEDISEG's deploy split (40 query pills): more enrolment photos help only partly.}
  \label{fig:robust}
\end{figure}

\begin{figure}[t]
  \centering
  \includegraphics[width=\linewidth]{fig13_identify_confidence.pdf}
  \caption{Calibrated confidence for identification. (a) Stated confidence against observed accuracy on held-out
    medicines. (b) Verdicts for catalogue pills and for pills not in the catalogue (MEDISEG).}
  \label{fig:confidence}
\end{figure}

\subsection{Comparison with the ePillID benchmark}
\label{sec:epillid}

For a like-for-like comparison, runs 5--7 follow ePillID's protocol and train on the references of all types
(Table~\ref{tab:epillid}, Figure~\ref{fig:epillid}). Our embedding alone (MAP 0.57--0.60) is in the range of the
paper's plain classifiers, not its best model. Re-ranking by imprint, which the paper's model did not use, raises MAP
to 0.77--0.78 and GAP@1 to 0.921, above the paper's 0.912, while MAP@1 stays lower (69.0 vs.\ 92.0\%). The paper's
ablations attribute its advantage to its training recipe (contrastive and hard-mined triplet losses with softmax and
ArcFace over side-specific classes, a 2,048-d embedding) rather than its backbone. We ruled out our crop style, the
64-d pooling, and local versus global pooling as causes.

\begin{table}[t]
  \caption{ePillID, official protocol (fold 4, two-sided). The best published model is a ResNet-152 with compact
    bilinear pooling and multi-head metric learning. ``+ imprint'' uses OCR information the published models did not
    use.}
  \label{tab:epillid}
  \centering
  \small
  \begin{tabular}{@{}lccccc@{}}
    \toprule
    & MAP & GAP & MAP@1 & GAP@1 & top-5 \\
    \midrule
    \citet{usuyama2020epillid}, best & \textbf{0.958} & \textbf{0.812} & \textbf{92.0} & 0.912 & -- \\
    \citet{usuyama2020epillid}, plain classifiers & .547--.727 & .241--.448 & 35.7--56.8 & .331--.551 & -- \\
    Ours, run 5 & 0.597 & 0.207 & 45.7 & 0.526 & 77.5 \\
    Ours, run 5 + imprint & 0.780 & 0.598 & 69.0 & \textbf{0.921} & 89.6 \\
    Ours, run 7 + imprint & 0.774 & 0.614 & 67.6 & 0.916 & 89.5 \\
    \bottomrule
  \end{tabular}
\end{table}

\begin{figure}[t]
  \centering
  \includegraphics[width=\linewidth]{fig10_epillid_benchmark.pdf}
  \caption{ePillID benchmark (official protocol): MAP@1 and GAP@1 of our runs with and without imprint re-ranking
    against the best published model.}
  \label{fig:epillid}
\end{figure}

\subsection{Error analysis and interpretation}
\label{sec:errors}

\paragraph{Imprints.} The reader finds a token of the pill's own imprint on about 47\% of validation photos.
Figure~\ref{fig:imprint} shows where it fails: debossed, small and round pills are hardest, and C3PI's own raters
marked many imprints unreadable.

\begin{figure}[t]
  \centering
  \includegraphics[width=\linewidth]{fig11_imprint_reading.pdf}
  \caption{Share of validation photos on which a token of the pill's own imprint is read (Florence-2-large, zero-shot),
    by imprint type, C3PI rater's judgement of visibility, pill size and shape.}
  \label{fig:imprint}
\end{figure}

\paragraph{Look-alikes.} The hardest errors are medicines that look the same (Figure~\ref{fig:lookalike}). With two
strengths of one ingredient in a gallery, one is named as the other 0.11\% of the time, or 0.58\% when shape and colour
also match. The imprint check settles retakes but cannot confirm strength, so strength is taken from the pharmacy
label. Same-manufacturer pairs that share an imprint prefix are confused 0.01\% of the time overall, but the worst
pairs share Teva's `TV', e.g.\ tramadol 50\,mg (TV;58) and carvedilol 12.5\,mg (TV;7295) at 5.8\%.

\begin{figure}[t]
  \centering
  \includegraphics[width=\linewidth]{fig16_lookalike_error_analysis.pdf}
  \caption{Look-alikes, with both medicines enrolled in one gallery. (a) Wrong-name rates by kind of pair, look only
    and with the imprint check. (b) The worst same-manufacturer pairs.}
  \label{fig:lookalike}
\end{figure}

\paragraph{What the model learned.} Photos of a medicine never seen in training form tight clusters
(Figure~\ref{fig:tsne}), and the calibrated threshold separates a photo's similarity to its own medicine from its
similarity to the closest of nine others (Figure~\ref{fig:sim}). Tracing the served model on single photos shows two
things. Capsule routing moves from equal weights to about 9\% CNN and 91\% ViT, and the refinement block's token
attention is nearly uniform, consistent with that block not being a bottleneck. Overall, the hybrid learns an
embedding that transfers to unseen medicines when it is trained with the right learning rates. Calibration gives
predictable behaviour on new medicines in familiar conditions but not across phones, and the imprint is evidence that
complements the look.

\begin{figure}[t]
  \centering
  \begin{minipage}[t]{0.49\linewidth}
    \centering
    \includegraphics[width=\linewidth]{fig14_similarity_distributions.pdf}
    \caption{Held-out photos: similarity to their own medicine's 3-photo prototype and to the closest of 9 other
      medicines, with the threshold $\tau$.}
    \label{fig:sim}
  \end{minipage}\hfill
  \begin{minipage}[t]{0.49\linewidth}
    \centering
    \includegraphics[width=\linewidth]{fig15_embedding_tsne.pdf}
    \caption{t-SNE~\citep{vandermaaten2008tsne} of the 4,005 held-out photos' embeddings; three medicines highlighted
      (photo counts in brackets).}
    \label{fig:tsne}
  \end{minipage}
\end{figure}

\paragraph{Negative results.} Table~\ref{tab:negative} lists what did not work. Two LoRA~\citep{hu2022lora} fine-tunes
of the imprint reader read more imprints on validation (44.0$\to$53.5\% and 45.0$\to$50.7\%). The first wrote
memorised imprints onto new pills; the second doubled wrong names in galleries (30$\to$61). Both were rejected, as were
look-alike-character matching and per-candidate imprint scoring.

\begin{table}[t]
  \caption{What did not work (validation).}
  \label{tab:negative}
  \centering
  \small
  \begin{tabular}{@{}p{0.34\linewidth}p{0.46\linewidth}p{0.12\linewidth}@{}}
    \toprule
    Tried & Result & Decision \\
    \midrule
    LoRA fine-tune of Florence-2-large, v1 (19,312 photos of training medicines) & imprints read 44.0$\to$53.5\%, but
      memorised imprints written onto new pills; identification and gallery check worse & rejected \\
    LoRA v2 (also learns to output nothing on unreadable photos) & own imprint read 45.0$\to$50.7\%; identification
      82.0$\to$83.0\%; gallery wrong names 30$\to$61 & rejected \\
    Look-alike characters (O/0, B/8) treated as equal & measured and dropped & rejected \\
    Scoring each candidate's imprint separately & measured and dropped & rejected \\
    Runs 6--7: stronger fine-tuning; benchmark-aligned training & better validation and MEDISEG, not ePillID & not
      served \\
    \bottomrule
  \end{tabular}
\end{table}

\subsection{Software 3.0 in coding and evaluation}
\label{sec:sw3-experiments}

The evaluation harness was written with the coding agent from natural-language specifications: the patient
simulation, the look-alike and same-manufacturer stress tests, the top-$k$ and confidence evaluations, the figures, and
the API's 92 tests. Tests and review caught real defects, among them a safety bug in the label reader that linked an
extended-release sticker to the immediate-release product; we fixed it by enforcing the release form. Our safeguards
against over-trusting generated code were procedural: settings were chosen on validation, each test set was scored
once, and every number in this paper is read from a file written by a script.
% [TO DO: the team decides how to evidence tool use (PROGRESS_LOG.md, git history); do not create mock logs.]
